\section{Introduction}
For a long time, genome mapping used a single linear reference sequence. That choice is convenient but biased: any sample allele absent from the reference is hard to see, and the bias falls unevenly across ancestries \cite{Cui2025, Secomandi2025}. Pangenome variation graphs were introduced to address this issue. Rather than one string, a pangenome graph encodes many haplotypes as paths through a graph whose nodes carry DNA sequences and whose ``bubbles'' make polymorphism explicit \cite{Garrison2024}. The Human Pangenome Reference Consortium (HPRC) draft, built from 47 diploid samples, is a good example for that \cite{Liao2023-short}.

This success creates a new need. Once you have pangenomes, for example an African pangenome and an East Asian pangenome, then come the questions:
\begin{enumerate}
    \item how do they differ?
    \item Which structural variants are shared, which are population-specific?
    \item and where exactly do the graphs diverge at base-pair resolution?
\end{enumerate}

Answering this requires aligning one graph to another graph. Surprisingly, no tool does this. In the current state of the art, the field has solved sequence-to-sequence alignment \cite{NEEDLEMAN1970443,SMITH1981195} and sequence-to-graph alignment \cite{Cui2025,santa_paper,segram_damla_onur}, but graph-to-graph alignment at the base-pair level remains open. In practice, comparing two pangenomes still means projecting both onto a common linear reference and comparing the projections, which reintroduces the very reference bias pangenomes exist to eliminate \cite{eizenga2020pangenome, Liao2023-short}. Another method is to analyze both pangenomes by creating some metrics and comparing those metrics against each other \cite{Hyun2022}.

In general, graph-to-graph alignment is NP-hard: given a subgraph-isomorphism instance $(H,G)$, label every node with a distinct sequence and set $\beta{=}1,\;\alpha{=}\gamma{=}0$; the optimal alignment score equals $|E(H)|$ iff $H$ embeds in $G$, so subgraph isomorphism \cite{garey1979computers} reduces to it in polynomial time. The key observation is that variation graphs are not arbitrary graphs. They decompose into a hierarchy of snarls, the generalised bubbles of Paten et al. \cite{paten2018superbubbles}, and this snarl tree has bounded tree-width in practice \cite{bounded-tree-width}.

We turn this observation into G2G, a system whose contribution is as much about how to compute the alignment efficiently as about the alignment itself. The workload is irregular: most snarls are tiny SNP bubbles, but a few structural-variant regions are enormous, so per-snarl work is heavy-tailed and a static work partition is badly load-imbalanced. G2G addresses this with three systems techniques: a cache-oblivious memory layout for the snarl tree, a minimiser index for anchor detection, and a work-stealing scheduler keyed on estimated subtree work, each motivated by direct profiling of the workload.

\paragraph{Contributions.}
\begin{itemize}
    \item \textbf{A new problem and a working system.} We define base-pair-resolution graph-to-graph alignment for variation graphs and present \textsc{G2G}, the first system to solve it. The delta graph it outputs is the graph-level analogue of a VCF: every allele present in one population graph but absent from the other.
    
    \item \textbf{A workload characterisation that drives the design.} Using the reference prototype we measure where time and memory go and show per-snarl work is heavy-tailed (Gini $>0.55$) and scale-invariant, motivating work stealing over static partitioning.
    
    \item \textbf{A parallel, cache-aware implementation.} A C++20 engine with a van~Emde~Boas snarl-tree layout, minimiser anchor index, exact LAPJV bipartite matching, and an OpenMP work-stealing runtime.
    
    \item \textbf{Evaluation on synthetic and real pangenomes,} including a reference-anchored diff that scales near-linearly past whole-genome size, and real African vs.\ East Asian (1000 Genomes) population graphs whose delta is significantly enriched near gnomAD-SV and ClinVar pathogenic sites ($p<10^{-3}$), demonstrating the delta lands on real structural and clinical variation.
\end{itemize}